\chapter{Counterparty Exposure}\label{ch:rc:counterparty-exposure}

When Lehman Brothers filed for bankruptcy on 15 September 2008 it was a party to around
930\,000 derivative transactions, by its own first count. Each counterparty now held
contracts with a defaulted firm: it could terminate them, net what it owed against what
it was owed, keep the collateral it held, and replace the hedges in a market that every
other counterparty was entering on the same side at the same time. By November most
contracts had been terminated; twenty months later, banks' claims for derivatives losses
exceeded 50 billion dollars. How much a bank stands to lose when a counterparty defaults
depends on what the contracts will be worth then, which is unknown today: the loss is an
option on a portfolio's future value. This chapter measures it. It opens Part~III, whose
chapters turn the exposure into prices (credit and debit valuation adjustments), into
funding, margin and capital costs, and into a desk that manages them.

\section{Exposure and its profiles}

\begin{definition}[Counterparty credit risk]\label{def:rc:counterparty-exposure:ccr}
\emph{Counterparty credit risk}\index{counterparty credit risk} is the risk that the other
party to a derivative defaults before the contract's final payment while the contract has
positive value to the surviving party. Unlike the credit risk of a loan, the amount at risk
is uncertain and can change sign.
\end{definition}

\begin{definition}[Counterparty exposure]\label{def:rc:counterparty-exposure:exposure}
The \emph{counterparty exposure}\index{counterparty exposure} at a future date $t$ is the
loss if the counterparty defaults then and nothing is recovered: $E_t = \max(V_t-C_t,0)$,
with $V_t$ the value of the contracts to the bank and $C_t$ the collateral it can keep.
\end{definition}

\begin{definition}[Exposure profiles]\label{def:rc:counterparty-exposure:profiles}
The \emph{expected exposure}\index{expected exposure} is $\mathrm{EE}(t) = \E[E_t]$; the
\emph{expected negative exposure}\index{expected negative exposure} is
$\mathrm{ENE}(t) = \E[\min(V_t-C_t,0)]$, the counterparty's expected exposure to the bank; the
\emph{potential future exposure}\index{potential future exposure} $\mathrm{PFE}_q(t)$ is the
$q$-quantile of $E_t$ (97.5\% here); the \emph{expected positive
exposure}\index{expected positive exposure} over a horizon $H$ is the time average
$\mathrm{EPE} = \frac1H\int_0^H\mathrm{EE}(t)\,dt$.
\end{definition}

EE and ENE price the adjustments of the next chapters; PFE sets credit limits (the most
the bank might lose with a given confidence); EPE is the one-number summary behind
regulatory exposure measures.

\begin{example}[Two trades with one counterparty]\label{ex:rc:counterparty-exposure:profiles}
A bank receives fixed at the par rate of 3.80\% on a USD~100 million ten-year swap
(chapter~1's SOFR curve, Hull--White with $\kappa = 3\%$ and a normal volatility of 90 basis
points), and with the same counterparty receives USD fixed at 3.80\% on USD~110 million and
pays EUR fixed at the ESTR par rate of 2.59\% on EUR~100 million for ten years, exchanging
the notionals at the end (EURUSD at 1.10, volatility 8\%). On 4\,000 paths simulated every
ten business days, the swap's EE peaks at USD~3.13 million after three years and its PFE
at USD~18.26 million after 3.9 years: rate risk grows with time and the remaining duration
shrinks. The cross-currency swap's EE rises to USD~7.09 million and its PFE to
USD~39.93 million just before maturity, driven by the final exchange
(\cref{fig:rc:counterparty-exposure:profiles}).
\end{example}

\begin{omfigure}
\begin{tikzpicture}
\begin{axis}[width=11.5cm,height=5.4cm,
  xlabel={time (years)},ylabel={USD million},
  tick label style={font=\small},label style={font=\small},
  xmin=0,xmax=10,ymin=0,ymax=42,grid=major,grid style={gray!25},
  legend columns=2,legend cell align=left,legend style={font=\footnotesize,at={(0.5,-0.25)},anchor=north,draw=none,fill=none,/tikz/every even column/.append style={column sep=8pt}},
  table/col sep=comma]
\addplot[omDef,very thick] table[x=t,y=swap_ee]{figdata/rates-credit-risk/17-counterparty-exposure/profiles.csv};
\addplot[omDef,thick,dashed] table[x=t,y=swap_pfe]{figdata/rates-credit-risk/17-counterparty-exposure/profiles.csv};
\addplot[omThm,very thick] table[x=t,y=xccy_ee]{figdata/rates-credit-risk/17-counterparty-exposure/profiles.csv};
\addplot[omThm,thick,dashed] table[x=t,y=xccy_pfe]{figdata/rates-credit-risk/17-counterparty-exposure/profiles.csv};
\legend{swap EE,swap PFE (97.5\%),cross-currency EE,cross-currency PFE (97.5\%)}
\end{axis}
\end{tikzpicture}

\omcaption{Exposure profiles of a ten-year interest rate swap (humped: uncertainty grows
while the remaining cash flows run off) and a ten-year cross-currency swap (rising to
maturity: the final exchange of notionals keeps the FX risk to the end). The small teeth
are the annual payments. Data: the chapter's tutorial.}\label{fig:rc:counterparty-exposure:profiles}
\end{omfigure}

\section{Netting and collateral}

\begin{definition}[Netting set, close-out netting]\label{def:rc:counterparty-exposure:netting}
A \emph{netting set}\index{netting set} is the group of trades with one counterparty under
one legally enforceable master agreement. \emph{Close-out netting}\index{close-out netting}
lets the surviving party terminate every trade of the netting set on default and set the
values off against each other into one net amount, so the exposure is
$\max(\sum_iV_i,0)$, not $\sum_i\max(V_i,0)$.
\end{definition}

\begin{dated}{2026-09}{Netting opinions}\label{dat:rc:counterparty-exposure:opinions}
Whether \omterm{def:rc:counterparty-exposure:netting}{close-out netting} holds is a question of each counterparty's insolvency law. ISDA has
published legal opinions on the enforceability of netting under its master agreements covering over 90
jurisdictions, and on collateral covering over 60, generally updated every year.
\end{dated}

\begin{example}[The netting benefit]\label{ex:rc:counterparty-exposure:netting}
Netted, the two trades of \cref{ex:rc:counterparty-exposure:profiles} have an EE that peaks
at USD~7.77 million after 3.9 years, where the two standalone EEs add to USD~9.24 million:
a 16\% saving (\cref{fig:rc:counterparty-exposure:netting}). The \omterm{def:rc:counterparty-exposure:netting}{netting set}'s ENE reaches
$-\text{USD}~21.0$ million: the bank pays the lower-rate currency and the final exchange is
at a forward rate above spot, so the trade drifts in the counterparty's favour.
\end{example}

\begin{omfigure}
\begin{tikzpicture}
\begin{axis}[width=11.5cm,height=5.4cm,
  xlabel={time (years)},ylabel={USD million},
  tick label style={font=\small},label style={font=\small},
  xmin=0,xmax=10,ymin=-25,ymax=12,grid=major,grid style={gray!25},
  legend columns=3,legend style={font=\footnotesize,at={(0.5,-0.25)},anchor=north,draw=none,fill=none,/tikz/every even column/.append style={column sep=8pt}},
  table/col sep=comma]
\addplot[omExo,thick,dashed] table[x=t,y=sum_ee]{figdata/rates-credit-risk/17-counterparty-exposure/netting.csv};
\addplot[omDef,very thick] table[x=t,y=net_ee]{figdata/rates-credit-risk/17-counterparty-exposure/netting.csv};
\addplot[omThm,very thick] table[x=t,y=net_ene]{figdata/rates-credit-risk/17-counterparty-exposure/netting.csv};
\legend{sum of standalone EEs,netted EE,netted ENE}
\end{axis}
\end{tikzpicture}

\omcaption{Netting the swap and the cross-currency swap: the netted EE lies below the sum of
the standalone EEs, and the netted ENE is larger in size than the EE because the \omterm{def:rc:counterparty-exposure:netting}{netting set}
drifts in the counterparty's favour. Data: the chapter's tutorial.}\label{fig:rc:counterparty-exposure:netting}
\end{omfigure}

The credit support annex of One Quant Book~2, chapter~10 makes the parties exchange
variation margin as the \omterm{def:rc:counterparty-exposure:netting}{netting set}'s value moves. Three terms set how much.

\begin{definition}[Collateral threshold, minimum transfer amount, independent amount]\label{def:rc:counterparty-exposure:csa}
The \emph{collateral threshold}\index{collateral threshold} $\mathrm{Th}$ is the exposure a
party accepts before calling for collateral: the collateral called is $V-\mathrm{Th}$ above it. The
\emph{minimum transfer amount}\index{minimum transfer amount} (MTA) is the smallest change
in collateral that is called. The \emph{independent amount}\index{independent amount} is
collateral posted regardless of the value, a buffer against moves during the time it takes
to close out.
\end{definition}

\section{Simulating exposure}

\begin{method}[Exposure by simulation]\label{met:rc:counterparty-exposure:sim}
(1) Choose a time grid, including the trades' payment dates. (2) Simulate the risk
factors jointly on the grid (here the Hull--White factor exactly, and EURUSD around its
forward). (3) Revalue every trade on every path and date, in closed form where possible, by
regression otherwise (the Longstaff--Schwartz method of One Quant Book~5). (4) Sum each
\omterm{def:rc:counterparty-exposure:netting}{netting set}; apply the collateral agreement path by path. (5) Average and take quantiles
across paths.
\end{method}

\omcode{code/firm/exposure/firm_exposure.py}{58}{78}{Scenario generation: the Hull--White factor on the grid with its exact drift and variance, and EURUSD around its forward.}\label{lst:rc:counterparty-exposure:sim}

The cost is paths times dates times trades times the cost of a valuation: millions of
valuations for one \omterm{def:rc:counterparty-exposure:netting}{netting set}, billions for a bank. The revaluation step is where
engines spend their time; the aggregation step (netting, collateral, profiles) runs over
very large arrays and is the part the firm writes in C++20 and Rust.

\section{The margin period of risk and the collateralised profile}

Collateral does not remove the exposure: when a counterparty defaults, the last collateral
it posted reflects the value at the last successful margin call, and the surviving party
needs time to notice the failure, dispute, terminate and re-hedge.

\begin{definition}[Margin period of risk]\label{def:rc:counterparty-exposure:mpor}
The exposure under a collateral agreement at $t$ is $\max(V_t-C_{t-\delta},0)$, where $\delta$ is the
margin period of risk of One Quant Book~1: the collateral is the one agreed $\delta$ earlier.
The exposure is the move of the \omterm{def:rc:counterparty-exposure:netting}{netting set}'s value over $\delta$, plus the threshold and the
\omterm{def:rc:counterparty-exposure:csa}{minimum transfer amount}.
\end{definition}

\omcode{code/firm/exposure/firm_exposure.py}{149}{165}{The collateral held under a two-way CSA, and the exposure that uses the collateral agreed one margin period of risk earlier.}\label{lst:rc:counterparty-exposure:csa}

\begin{omfigure}
\begin{tikzpicture}[x=1.05cm,y=1cm,font=\small]
\draw[->,thick] (0,0) -- (10.8,0) node[right] {$t$};
\foreach \x in {1,2,3,4} {\draw[omDef,thick] (\x,0.12) -- (\x,-0.12);}
\node[omDef,below,align=center] at (2.5,-0.2) {daily margin calls met};
\draw[omThm,very thick] (5,0.3) -- (5,-0.3);
\node[omThm,above,align=center] at (5,0.35) {last call met\\$C = V_{t-\delta}$};
\draw[omThm,very thick,dashed] (7,0.3) -- (7,-0.3);
\node[omThm,above] at (7,0.35) {default};
\draw[omThm,very thick] (9.6,0.3) -- (9.6,-0.3);
\node[omThm,above,align=center] at (9.6,0.35) {close-out\\$V_t$};
\draw[<->,omExo,thick] (5,-0.75) -- (9.6,-0.75);
\node[omExo,below,align=center] at (7.3,-0.8) {margin period of risk $\delta$: no collateral moves;\\exposure $= \max(V_t - C,\,0)$};
\end{tikzpicture}

\omcaption{The margin period of risk: the last collateral received matches the value at
the last successful call; between it and the close-out the counterparty posts nothing, and
the value keeps moving. Schematic.}\label{fig:rc:counterparty-exposure:mpor}
\end{omfigure}

\begin{example}[Collateralising the netting set]\label{ex:rc:counterparty-exposure:csa}
Under a two-way CSA with zero threshold and a \omterm{def:rc:counterparty-exposure:csa}{minimum transfer amount} of USD~500\,000, the
\omterm{def:rc:counterparty-exposure:netting}{netting set}'s EE falls to a nearly flat USD~1.30 million at its peak and its PFE to
USD~6.45 million with a margin period of risk of ten business days; the EPE over the life
falls from USD~6.44 million to 0.95 million, 15\% of the uncollateralised figure. With
twenty days the peak EE is USD~1.83 million, $\sqrt2$ times higher; with a threshold of
USD~10 million it is USD~3.81 million (\cref{fig:rc:counterparty-exposure:collat}).
\end{example}

\begin{omfigure}
\begin{tikzpicture}
\begin{axis}[width=11.5cm,height=5.4cm,
  xlabel={time (years)},ylabel={expected exposure (USD million)},
  tick label style={font=\small},label style={font=\small},
  xmin=0,xmax=10,ymin=0,ymax=9,grid=major,grid style={gray!25},
  legend columns=2,legend cell align=left,legend style={font=\footnotesize,at={(0.5,-0.25)},anchor=north,draw=none,fill=none,/tikz/every even column/.append style={column sep=8pt}},
  table/col sep=comma]
\addplot[omExo,very thick] table[x=t,y=none]{figdata/rates-credit-risk/17-counterparty-exposure/collat.csv};
\addplot[omMeth,very thick] table[x=t,y=th10]{figdata/rates-credit-risk/17-counterparty-exposure/collat.csv};
\addplot[omThm,thick] table[x=t,y=csa20]{figdata/rates-credit-risk/17-counterparty-exposure/collat.csv};
\addplot[omDef,thick] table[x=t,y=csa10]{figdata/rates-credit-risk/17-counterparty-exposure/collat.csv};
\legend{no CSA,threshold USD~10 million (10 days),zero threshold (20 days),zero threshold (10 days)}
\end{axis}
\end{tikzpicture}

\omcaption{\omterm{def:rc:counterparty-exposure:profiles}{Expected exposure} of the \omterm{def:rc:counterparty-exposure:netting}{netting set} without collateral and under three CSAs. A
zero-threshold CSA flattens the profile to the value moves over the margin period of risk;
doubling that period multiplies them by about $\sqrt2$; a threshold adds exposure up to its
size. The teeth at each anniversary are the annual payments, which the lagged collateral
does not yet reflect. Data: the chapter's tutorial.}\label{fig:rc:counterparty-exposure:collat}
\end{omfigure}

\begin{dated}{2026-09}{Margin period of risk floors}\label{dat:rc:counterparty-exposure:mpor}
The Basel Committee's standardised approach for \omterm{def:rc:counterparty-exposure:ccr}{counterparty credit risk} (March 2014)
sets a minimum margin period of risk of ten business days for bilateral derivatives with
daily margin, five for cleared client trades and twenty for \omterm{def:rc:counterparty-exposure:netting}{netting sets} of 5\,000 or more
transactions not with a central counterparty, doubled for \omterm{def:rc:counterparty-exposure:netting}{netting sets} with outstanding
disputes; with margin every $N$ days it is $10+N-1$ days.
\end{dated}

\section{Wrong-way risk}

\begin{definition}[Wrong-way risk]\label{def:rc:counterparty-exposure:wwr}
\emph{Wrong-way risk}\index{wrong-way risk} is a positive dependence between a
counterparty's \omterm{def:rc:reduced-form-credit:pd}{probability of default} and the bank's exposure to it: the exposure is
largest when default is likeliest. It is \emph{specific} when the dependence is built into
the trade (protection bought from a counterparty on itself or on its own country) and
\emph{general} when it comes through common market factors.
\end{definition}

\begin{example}[A counterparty that weakens with the euro]\label{ex:rc:counterparty-exposure:wwr}
Let the counterparty of the cross-currency swap have a hazard rate of $2\%\times(X_t/F_X(0,t))^{-5}$:
a 10\% fall of the euro below its forward raises its hazard by about 70\%. The bank's
exposure is largest when the euro has fallen. Weighting the paths by the \omterm{def:rc:reduced-form-credit:pd}{probability of
default} in each period, the \omterm{def:rc:counterparty-exposure:profiles}{expected exposure} conditional on default is 2.52 times the
unconditional EE at its peak, and 2.22 times on average over the life
(\cref{fig:rc:counterparty-exposure:wwr}).
\end{example}

\begin{omfigure}
\begin{tikzpicture}
\begin{axis}[width=11.5cm,height=5.4cm,
  xlabel={time (years)},ylabel={USD million},
  tick label style={font=\small},label style={font=\small},
  xmin=0,xmax=10,ymin=0,ymax=20,grid=major,grid style={gray!25},
  legend columns=2,legend cell align=left,legend style={font=\footnotesize,at={(0.5,-0.25)},anchor=north,draw=none,fill=none,/tikz/every even column/.append style={column sep=8pt}},
  table/col sep=comma]
\addplot[omDef,very thick] table[x=t,y=ee]{figdata/rates-credit-risk/17-counterparty-exposure/wwr.csv};
\addplot[omThm,very thick,dashed] table[x=t,y=cee]{figdata/rates-credit-risk/17-counterparty-exposure/wwr.csv};
\legend{EE,EE given default (wrong-way)}
\end{axis}
\end{tikzpicture}

\omcaption{\omterm{def:rc:counterparty-exposure:profiles}{Expected exposure} of the cross-currency swap, unconditionally and conditional on
the counterparty defaulting at each date, when the counterparty's hazard rises as the euro
falls. The exposures that matter are those in the scenarios where the counterparty
defaults. Data: the chapter's tutorial.}\label{fig:rc:counterparty-exposure:wwr}
\end{omfigure}

\begin{dated}{2026-09}{Lehman's derivatives}\label{dat:rc:counterparty-exposure:lehman}
According to the Financial Crisis Inquiry Commission, Lehman first asserted around
930\,000 derivative transactions at its bankruptcy; most had been terminated by
13 November 2008, 18\,000 were still outstanding in January 2009, and by May 2010 banks had
filed more than 50 billion dollars of claims for losses on derivatives with Lehman.
\end{dated}

\section{Tutorial: exposure of a netting set}\label{tut:rc:counterparty-exposure:tutorial}

\begin{tutorial}
\textbf{Goal.} Simulate the exposure of a swap, a cross-currency swap and their \omterm{def:rc:counterparty-exposure:netting}{netting set},
with and without variation margin. \textbf{End state:} the numbers of
\cref{ex:rc:counterparty-exposure:profiles,ex:rc:counterparty-exposure:netting,ex:rc:counterparty-exposure:csa,ex:rc:counterparty-exposure:wwr}
and their four charts.
\begin{tutsteps}
\item \textbf{Scenarios}: \texttt{simulate(MKT, 10, 26, 4000)}; check that EURUSD is a
martingale around its forward.
\item \textbf{Values}: \texttt{Swap.values(sc)}, \texttt{XccySwap.values(sc)}; both
are worth zero at the start.
\item \textbf{Profiles}: \texttt{aggregate([swap])}, the \omterm{def:rc:counterparty-exposure:netting}{netting set}, and the three CSAs.
\item \textbf{Wrong way}: \texttt{wrong\_way()}; \texttt{fig\_rc\_exposure.py} writes the charts.
\end{tutsteps}
\textbf{What to change next.} Make the bank pay fixed on the swap and watch the netting
benefit change; build the C++20 kernel (\texttt{cpp/}) and time it against the Python
aggregation on 100\,000 paths.
\end{tutorial}

\section{Build: the exposure engine}\label{bld:rc:counterparty-exposure:exposure}

\begin{build}
\bfield{Purpose} The firm's exposure engine: profiles of every \omterm{def:rc:counterparty-exposure:netting}{netting set} for credit limits
(PFE), for the valuation adjustments of chapters~18 to~20 (EE, ENE), and for the risk
engine of chapter~29.
\bfield{Interface} \texttt{Market}, \texttt{simulate(market, horizon, steps\_per\_year,
paths, seed)} returning \texttt{Scenarios} (\texttt{usd\_bond}, \texttt{eur\_bond});
trades \texttt{Swap}, \texttt{XccySwap}, \texttt{FxForward} with \texttt{values(sc)};
\texttt{collateral}, \texttt{exposure}, \texttt{profiles}, \texttt{epe},
\texttt{aggregate}, \texttt{conditional\_ee}. C++20 \texttt{firm::exposure::aggregate}
(\texttt{cpp/firm\_exposure.hpp}) and Rust \texttt{firm\_exposure::aggregate}.
\bfield{Rules} Risk-neutral scenarios; one grid step is the margin period of risk (ten
business days); the final grid date, when all trades have matured, is excluded from
profiles; PFE is the lower 97.5\% quantile.
\bfield{Acceptance tests} \texttt{code/firm/exposure/tests/}, \texttt{cpp/} and
\texttt{rust/}: the three implementations reproduce one fixture to $10^{-9}$; par trades are
worth zero; FX and deflated bonds are martingales; netting and collateral reduce the peak
EE, and a longer margin period of risk raises it.
\bfield{Stretch} Cash flows during the margin period of risk; stochastic EUR rates and
rate--FX correlation; regression revaluation of callable trades; the C++ kernel called
from Python.
\end{build}

\begin{omsources}
\item Basel Committee on Banking Supervision, \emph{The standardised approach for measuring
counterparty credit risk exposures}, March 2014.
\item Financial Crisis Inquiry Commission, \emph{The Financial Crisis Inquiry Report},
2011, notes to chapter~20.
\item J.~Gregory, \emph{The xVA Challenge: Counterparty Risk, Funding, Collateral, Capital and Initial
Margin}, 4th edition, Wiley, 2020.
\end{omsources}

\section{Exercises}

\begin{exercise}[$\star$]\label{exo:rc:counterparty-exposure:1}
Why is the \omterm{def:rc:counterparty-exposure:profiles}{expected exposure} of an interest rate swap humped, and that of a cross-currency
swap with final exchange increasing?
\end{exercise}

\begin{exercise}[$\star$]\label{exo:rc:counterparty-exposure:2}
Two trades with one counterparty are worth $+10$ and $-7$ on a path. What is the exposure
with and without \omterm{def:rc:counterparty-exposure:netting}{close-out netting}?
\end{exercise}

\begin{exercise}[$\star$]\label{exo:rc:counterparty-exposure:3}
A CSA has a threshold of USD~10 million and an MTA of USD~0.5 million; the \omterm{def:rc:counterparty-exposure:netting}{netting set} is
worth USD~14.3 million and the bank holds USD~4.0 million. How much does it call?
\end{exercise}

\begin{exercise}[$\star\star$]\label{exo:rc:counterparty-exposure:4}
Why does doubling the margin period of risk multiply the collateralised EE by about
$\sqrt2$?
\end{exercise}

\begin{exercise}[$\star\star$]\label{exo:rc:counterparty-exposure:5}
The netted ENE of \cref{ex:rc:counterparty-exposure:netting} is larger in size than the EE.
What does that mean for the counterparty, and for the bank's DVA (chapter~18)?
\end{exercise}

\begin{exercise}[$\star\star$]\label{exo:rc:counterparty-exposure:6}
Give two examples of specific \omterm{def:rc:counterparty-exposure:wwr}{wrong-way risk} and one of right-way risk.
\end{exercise}

\begin{exercise}[$\star\star\star$]\label{exo:rc:counterparty-exposure:7}
\emph{Coding.} Compute the peak EE of the \omterm{def:rc:counterparty-exposure:netting}{netting set} collateralised with a twenty-day
margin period of risk and check the ratio to the ten-day figure.
\end{exercise}

\begin{exercise}[$\star\star\star$]\label{exo:rc:counterparty-exposure:8}
\emph{Find the flaw.} ``The counterparty posts daily variation margin with zero threshold,
so our exposure to it is zero and it needs no credit limit.''
\end{exercise}

\section{Problem: The Uncollateralised Corporate}

\begin{problem}[{Weekend problem --- signing a CSA}]\label{pb:rc:counterparty-exposure:1}
A corporate client pays fixed to the bank on a USD~50 million five-year swap (at the par rate)
and buys EUR~20 million one year forward from the bank (at the forward rate). It has no CSA;
the bank proposes one with zero threshold and a \omterm{def:rc:counterparty-exposure:csa}{minimum transfer amount} of USD~250\,000, or
one with a threshold of USD~5 million.

\medskip\noindent\textbf{Part I --- Without a CSA.}
\begin{enumerate}
\item Give the swap's par rate and the forward's strike.
\item Give the \omterm{def:rc:counterparty-exposure:netting}{netting set}'s peak PFE and when it occurs.
\item Give its EPE over one year and over its life.
\item Which trade dominates the first year, and why?
\item Why does a corporate often refuse to sign a CSA?
\end{enumerate}

\medskip\noindent\textbf{Part II --- With a CSA.}
\begin{enumerate}[resume]
\item Give the peak PFE and the EPEs under the zero-threshold CSA.
\item Why does the collateralised PFE peak at one year exactly?
\item Give the peak PFE and the life EPE with the USD~5 million threshold.
\item Why does the USD~5 million threshold change almost nothing here?
\item What would the corporate need in order to post variation margin?
\end{enumerate}

\medskip\noindent\textbf{Part III --- Limits and capital.}
\begin{enumerate}[resume]
\item The bank's limit for the client is a PFE of USD~4 million. Is the client within it?
\item How does the limit change the bank's willingness to trade more?
\item How would the bank reduce the settlement-date PFE?
\item What happens to the exposure if the euro rises 10\% in the first month?
\item Which party is exposed to whom on the FX forward if the euro falls?
\end{enumerate}

\medskip\noindent\textbf{Part IV --- Judgement.}
\begin{enumerate}[resume]
\item Why do banks still trade uncollateralised with corporates?
\item How would a margin period of risk of twenty days change the answers?
\item What does the bank give up by accepting a threshold?
\item State the \emph{named result}: the peak PFE and the life EPE before and after the
zero-threshold CSA.
\item In one sentence: what does a CSA turn a \omterm{def:rc:counterparty-exposure:exposure}{counterparty exposure} into?
\end{enumerate}
\end{problem}

\section{Interview questions}

\begin{interviewq}[$\star$ \iqroles{risk, bank}]\label{iq:rc:counterparty-exposure:1}
Define EE, PFE and EPE. Which is used for what?
\end{interviewq}

\begin{interviewq}[$\star\star$ \iqroles{developer, risk}]\label{iq:rc:counterparty-exposure:2}
Design an exposure simulation for a bank's whole derivative book. Where does the time go?
\end{interviewq}

\begin{interviewq}[$\star\star$ \iqroles{risk}]\label{iq:rc:counterparty-exposure:3}
What is the margin period of risk and why is it not one day for a daily margined \omterm{def:rc:counterparty-exposure:netting}{netting
set}?
\end{interviewq}

\begin{interviewq}[$\star\star$ \iqroles{trader, risk}]\label{iq:rc:counterparty-exposure:4}
Give an example of \omterm{def:rc:counterparty-exposure:wwr}{wrong-way risk} and how you would measure it.
\end{interviewq}

\begin{interviewq}[$\star\star\star$ \iqroles{researcher, developer}]\label{iq:rc:counterparty-exposure:5}
How do you compute the exposure of a \omterm{def:rc:bermudans-and-callables:bermudan}{Bermudan swaption} inside a simulation?
\end{interviewq}

\begin{interviewq}[$\star\star\star$ \iqroles{risk, bank}]\label{iq:rc:counterparty-exposure:6}
What happened to Lehman's derivative counterparties in the weeks after 15 September 2008?
\end{interviewq}
